\documentclass[./eval.tex]{subfiles}
\newcommand{\head}{\noindent Máquinas Eléctricas II |Evaluación| Apellidos y Nombres.
\rule{10cm}{0.1mm}}
\newcommand{\separador}{\noindent\rule{\paperwidth}{0.1mm}\newline}

\begin{document}
\head

 \begin{enumerate}
 
 \item Determinar el numero de pares de polos que posee un generador que gira a una velocidad de $n=4500rpm$ y genera una frecuencia de 300Hz.

    \item Se conectan en triángulo bobinas de $R=10\Omega$ y $X_L=10\Omega$ a una red trifásica de $V_C=400V, f=50Hz$. Calcular $I_C,I_S, cos(\varphi), P,Q,S$.
 

    \item Un motor trifásico de $P=3990W$ y $cos(\varphi)=0,65$ se conecta a una red de $V_C=380V, f=50hz$, se trata de averiguar la corriente de línea y de cada fase del motor cuando está conectado en triángulo, así como su potencia activa $P$ reactiva $Q$ y aparente $S$. Si cada una de las fases del motor se lo puede considerar como una inductancia $Z$ que posee una reactancia inductiva $X_L$ y una resistencia ohmica $R$. Determina el valor de las mismas.
   
    \separador
\end{enumerate}
\end{document} 